\documentclass[aspectratio=169]{beamer}
\input{header}
\coursecode{JKSSB}

\title{Internet, WWW \& Browser}
\subtitle{Lesson 34 --- The Web and Its Access Tools}
\author{JKSSB Computer Awareness}
\date{\today}

\begin{document}
\frame{\titlepage}

\begin{frame}{What the Internet Is}
  \begin{itemize}
    \item The \key{Internet} is a global network of connected computers and
      networks.
    \item It is a \key{network of networks}: smaller networks joined together
      so any connected device can reach any other.
    \item It carries many services --- the Web, e-mail, file transfer, and
      more.
    \item The Internet is the \key{physical infrastructure}: cables, routers,
      servers, and the devices connected to them.
  \end{itemize}
  \begin{block}{The one idea}
    The Internet is the network itself; the services that run on it, such as
    the Web and e-mail, are not the same thing as the Internet.
  \end{block}
\end{frame}

\begin{frame}{Internet vs WWW}
  \begin{center}
    \begin{tabular}{lll}
      \toprule
      \textbf{Point} & \textbf{Internet} & \textbf{WWW (World Wide Web)} \\
      \midrule
      What it is & Global network of networks & A service running on the Internet \\
      Content & Carries all services & Linked web pages and sites \\
      Access & Via an ISP connection & Via a browser over the Internet \\
      \bottomrule
    \end{tabular}
  \end{center}
  \vspace{0.35cm}
  \begin{alertblock}{Trap alert}
    The Internet and the Web are \key{not} the same. The \key{WWW --- World
    Wide Web} --- is one service that uses the Internet.
  \end{alertblock}
\end{frame}

\begin{frame}{Website vs Webpage vs Homepage}
  \begin{itemize}
    \item A \key{webpage} is a single document on the Web, shown in a browser.
    \item A \key{website} is a collection of related webpages under one
      address (domain).
    \item The \key{homepage} is the main or starting page of a website.
    \item Order of size: \key{homepage} (one page) $\rightarrow$ \key{webpage}
      (any one page) $\rightarrow$ \key{website} (the whole collection).
  \end{itemize}
  \begin{block}{Keep the three straight}
    One page is a webpage; the whole collection is a website; the front door
    of that collection is the homepage.
  \end{block}
\end{frame}

\begin{frame}{What a Browser Is}
  \begin{itemize}
    \item A \key{browser} is software used to open and view websites and
      webpages.
    \item It sends a request for a page and displays the page that comes back.
    \item Common browsers: \key{Google Chrome}, \key{Mozilla Firefox},
      \key{Microsoft Edge}, \key{Apple Safari}.
    \item A browser is \key{software on the user's device} --- it is not the
      Internet itself.
  \end{itemize}
  \begin{block}{Office desktop example}
    On the office desktop, double-clicking the Chrome icon opens a browser;
    typing an address then loads that site's homepage.
  \end{block}
\end{frame}

\begin{frame}{What a Search Engine Is}
  \begin{itemize}
    \item A \key{search engine} is a website (service) that finds other
      webpages matching typed keywords.
    \item Examples: \key{Google}, \key{Bing}, \key{Yahoo}.
    \item It is \key{reached through} a browser --- the browser is the tool,
      the search engine is the site visited with it.
    \item \muted{Searching needs both: a browser open on the device, and a
      search-engine site loaded inside it.}
  \end{itemize}
\end{frame}

\begin{frame}{Browser vs Search Engine}
  \begin{center}
    \begin{tabular}{lll}
      \toprule
      \textbf{Point} & \textbf{Browser} & \textbf{Search engine} \\
      \midrule
      What it is & Software (application) & A website / service \\
      Role & Opens and displays pages & Finds pages for keywords \\
      Examples & Chrome, Firefox, Edge, Safari & Google, Bing, Yahoo \\
      \bottomrule
    \end{tabular}
  \end{center}
  \vspace{0.35cm}
  \begin{alertblock}{Trap alert}
    Chrome is a \key{browser}; Google is a \key{search engine}. Swapping the
    two is a common wrong answer.
  \end{alertblock}
\end{frame}

\begin{frame}{Client}
  \begin{itemize}
    \item A \key{client} is the device or program that requests a service or
      resource.
    \item The office PC's browser asking for a webpage is acting as a client.
    \item A client \key{asks}; it does not supply the page itself.
    \item \muted{Client and server always appear as a pair: one requests, the
      other serves.}
  \end{itemize}
  \begin{block}{The one idea}
    Client = the requester; server = the supplier.
  \end{block}
\end{frame}

\begin{frame}{Server}
  \begin{itemize}
    \item A \key{server} is a computer (or program) that supplies services or
      resources to clients.
    \item A \key{web server} stores websites and sends pages to browsers that
      ask for them.
    \item One server can serve \key{many clients} at the same time.
    \item Servers stay \key{switched on and connected} so clients can reach
      them at any time.
  \end{itemize}
\end{frame}

\begin{frame}{Client and Server Together}
  \begin{itemize}
    \item Step 1: the \key{client} (browser) sends a request for a page.
    \item Step 2: the \key{server} receives the request and sends the page
      back.
    \item Step 3: the browser \key{displays} the received page.
    \item This request-and-response pair is the \key{client-server} model of
      the Web.
  \end{itemize}
  \begin{exampleblock}{Worked example}
    Typing a site address and pressing Enter sends a request; the site's
    server returns the homepage, which the browser shows.
  \end{exampleblock}
\end{frame}

\begin{frame}{ISP: Internet Service Provider}
  \begin{itemize}
    \item \key{ISP} stands for Internet Service Provider.
    \item It is the company or organisation that \key{provides the Internet
      connection} to homes and offices.
    \item Without an ISP connection, the device cannot reach the Internet at
      all.
    \item The ISP link comes \key{before} everything else: no ISP, no
      browsing, no search, no e-mail.
  \end{itemize}
  \begin{alertblock}{Exam fact (PYQ-34)}
    ISP = \key{Internet Service Provider} (WO 2026 Q74). Do not expand it as
    anything else.
  \end{alertblock}
\end{frame}

\begin{frame}{Hyperlink}
  \begin{itemize}
    \item A \key{hyperlink} (link) is clickable text or an image that jumps to
      another page or section.
    \item Clicking it loads the \key{linked webpage}, which may sit on the
      same site or a different one.
    \item Links are what make the Web a ``web'' --- pages joined to pages.
    \item \muted{A hyperlink points the way; the browser follows it when
      clicked.}
  \end{itemize}
  \begin{block}{Keep the pair straight}
    A hyperlink \key{points} to a page; a browser \key{opens} it.
  \end{block}
\end{frame}

\begin{frame}{Full-Form Ladder of This Lesson}
  \begin{center}
    \begin{tabular}{lll}
      \toprule
      \textbf{Short form} & \textbf{Full form} & \textbf{One-line role} \\
      \midrule
      ISP & Internet Service Provider & Gives the Internet connection \\
      WWW & World Wide Web & Linked pages and sites \\
      URL & Uniform Resource Locator & Address of a page or resource \\
      HTTP & HyperText Transfer Protocol & Rules for requesting pages \\
      HTML & HyperText Markup Language & Language webpages are written in \\
      \bottomrule
    \end{tabular}
  \end{center}
  \vspace{0.25cm}
  \muted{Give the full form on first mention; direct-recall items ask the
  expansion exactly.}
\end{frame}

\begin{frame}{URL and Homepage Address}
  \begin{itemize}
    \item A \key{URL --- Uniform Resource Locator} --- is the address of a
      webpage or resource on the Web.
    \item Typing the site's URL in the browser's \key{address bar} loads its
      homepage.
    \item Each webpage of a website has \key{its own URL}; the homepage URL is
      the site's base address.
    \item \muted{The URL names the page; the hyperlink jumps to it.}
  \end{itemize}
\end{frame}

\begin{frame}{Trap Alert: Web Facts to Keep Straight}
  \begin{itemize}
    \item Internet $\neq$ WWW: the Web is \key{one service} on the Internet.
    \item Browser $\neq$ search engine: Chrome is a \key{browser}; Google is
      a \key{search engine}.
    \item Website $\neq$ webpage: the site is the \key{collection}; a
      webpage is \key{one document}; the homepage is the \key{first page}.
    \item Client \key{requests}; server \key{supplies}. (PYQ-34: ISP =
      Internet Service Provider.)
    \item A hyperlink \key{points}; only clicking (via the browser) follows
      it.
  \end{itemize}
\end{frame}

\begin{frame}{Lesson 34: Recall}
  \begin{itemize}
    \item The Internet is the global network of networks; the WWW --- World
      Wide Web --- is a service running on it.
    \item A webpage is one document; a website is the collection; the
      homepage is the main page.
    \item A browser (Chrome, Firefox, Edge, Safari) opens pages; a search
      engine (Google, Bing, Yahoo) finds pages.
    \item The client requests and the server supplies; the ISP --- Internet
      Service Provider --- gives the connection.
    \item A hyperlink jumps to another page; a URL --- Uniform Resource
      Locator --- is a page's address.
    \item Full forms: ISP, WWW, URL, HTTP, HTML --- learn each expansion
      exactly.
  \end{itemize}
\end{frame}
\end{document}
